% ---------------------------------------------------------------------------
% PROPOSED REPLACEMENT for tab:ccao_baseline_complementary (paper/paper_v20.tex, l.2359-2400).
% Label preserved.  Source: analysis/linear_baseline_reproduction_v1/tables/
% lb_table2_complementary_baseline.csv
% ---------------------------------------------------------------------------
\begin{table}[!htbp]
\centering
\scriptsize
\setlength{\tabcolsep}{2.4pt}
\renewcommand{\arraystretch}{1.06}
\newcommand{\baselinecompmetric}[1]{\hspace*{0.6em}#1}
\caption{Complementary baseline diagnostics for the ordinary linear benchmark and
\emph{Ordinary LightGBM (standard raw-label native)}: valuation level, horizontal uniformity,
and mechanism/residual structure. Descriptive levels for two unpenalized specifications, not a
recommendation between them.}
\label{tab:ccao_baseline_complementary}
\begin{tabularx}{\textwidth}{@{} >{\raggedright\arraybackslash}p{2.80cm} >{\centering\arraybackslash}X >{\centering\arraybackslash}X >{\centering\arraybackslash}X >{\centering\arraybackslash}X @{}}
\toprule
\textbf{Measure} & \multicolumn{2}{c}{\textbf{Held-out evaluation}} & \multicolumn{2}{c}{\textbf{2025 forward evaluation}} \\
\cmidrule(lr){2-3}\cmidrule(l){4-5}
& \textbf{Linear} & \textbf{\shortstack{Light\\GBM}} & \textbf{Linear} & \textbf{\shortstack{Light\\GBM}} \\
\midrule
\multicolumn{5}{@{}l}{\textbf{Valuation level}} \\
\cmidrule(r){1-1}
\baselinecompmetric{Median ratio $m_S$} & 0.969 & 0.929 & 1.020 & 0.950 \\
\baselinecompmetric{Mean ratio $\bar r_S$} & 1.020 & 0.989 & 1.081 & 1.015 \\
\baselinecompmetric{Weighted mean $\bar r_{W,S}$} & 0.979 & 0.924 & 1.027 & 0.941 \\
\addlinespace[5pt]
\multicolumn{5}{@{}l}{\textbf{Horizontal uniformity}} \\
\cmidrule(r){1-1}
\baselinecompmetric{COD} & 24.7\% & 21.6\% & 24.3\% & 21.3\% \\
\baselinecompmetric{COV} & 45.2\% & 39.7\% & 42.1\% & 37.2\% \\
\addlinespace[5pt]
\multicolumn{5}{@{}l}{\textbf{Mechanism and residual structure}} \\
\cmidrule(r){1-1}
\baselinecompmetric{$\beta_{\log}$} & -0.092 & -0.150 & -0.109 & -0.164 \\
\baselinecompmetric{$\Delta_{\mathrm{NL}}$} & 0.131 & 0.119 & 0.122 & 0.121 \\
\baselinecompmetric{$\operatorname{dCor}(e,y)$} & 0.250 & 0.387 & 0.269 & 0.422 \\
\bottomrule
\end{tabularx}
\vspace{1mm}
\begin{minipage}{\textwidth}
\scriptsize
\emph{Notes.}
Median, mean, weighted-mean ratio, COD, and COV are assessor-facing valuation-level or
horizontal-uniformity diagnostics defined in Section~\ref{subsec:metrics}. The first-order
mechanism slope $\beta_{\log}$, correlation-ratio nonlinearity gap $\Delta_{\mathrm{NL}}$, and
distance correlation are defined in Eqs.~\eqref{eq:log_ratio_slope},
\eqref{eq:nonlinearity_gap}, and~\eqref{eq:dcor_diagnostic}; the last two carry no
assessor-standard reference range and are residual-structure diagnostics rather than equity
measures.
Both columns' COD exceeds the $[5,15]$ range the adopted Standard on Ratio Studies gives for
single-family residential property (Table~\ref{tab:assessment_metrics_summary}), and COD rises
further along the penalty paths (Appendix Table~\ref{tab:path_anchor_complementary}). No
compliance claim is implied anywhere in this paper. This range comparison is a descriptive
interpretation of the reported values rather than a formal inferential finding.
The PRB and VEI value proxy is model-specific by construction: it is
$0.5\,P_i+0.5\,\widehat P_i/m_S$, so each column is read against a value axis that embeds that
column's own valuations and its own median ratio. The two columns are therefore not ranked on a
shared value axis.
COD and COV are reported in percent; COV is stored as a fraction in the retained artifact and
converted to percent for display.
\end{minipage}
\end{table}
